\normalfalse \difficiletrue \tdifficilefalse
\correctiontrue

%\UPSTIidClasse{11} % 11 sup, 12 spé
%\newcommand{\UPSTIidClasse}{12}

\exer{Pompe à piston axial  $\star$ \label{GEO:03:C2:06:11}}
\setcounter{question}{0}\marginnote{\xpComp{GEO}{03}}%\UPSTIcompetence[2]{C2-06}
\index{Compétence C2-06}\index{Compétence GEO-03}
\index{Pompe à pistons radiaux}
\index{Arbre à cames}
\ifcorrection
\else
\marginnote{\textbf{Pas de corrigé pour cet exercice.}}
\fi

\ifprof
\else
Soit le mécanisme suivant. On a $\vect{AB}=e\vect{i_1}$ et $\vect{BI}=R\vect{j_0}$ et $\vect{AC}=\lambda(t)\vect{j_0}$. De plus,
$e=\SI{10}{mm}$ et $R=\SI{20}{mm}$. Le contact entre \textbf{1} et \textbf{2} en $B$ est maintenu en permanence par un ressort suffisamment raide (non représenté) positionné entre \textbf{0} et \textbf{2}.
\begin{marginfigure}
\includegraphics[width=\linewidth]{11_01}
\end{marginfigure}
\fi

\question{Tracer le graphe des liaisons.}

\ifprof
\begin{marginfigure}
\includegraphics[width=\linewidth]{11_01_c}
\end{marginfigure}
\fi

\question{Exprimer $\lambda(t)$ en fonction de $\theta(t)$.}

\ifprof
En écrivant la fermeture géométrique, on a
$\vect{AB}+\vect{BI}+\vect{IC}+\vect{CA}=\vect{0}$.

On a donc
$e\vi{1}+R\vj{0}+\mu\vi{0}-\lambda(t)\vj{0}=\vect{0}$.
En projetant sur $\vj{0}$ :
\[
e\sin\theta(t)+R-\lambda(t)=0.
\]
Ainsi,
\[
\lambda(t)=e\sin\theta(t)+R.
\]
\fi

\question{Exprimer $\dot{\lambda}(t)$ en fonction de $\dot{\theta}(t)$.}
\ifprof
En dérivant l'expression précédente :
\[
\dot{\lambda}(t)=e\dot{\theta}(t)\cos\theta(t).
\]
\fi

\question{On note $S$ la section du piston \textbf{2}. Exprimer le débit instantané de la pompe.}
\ifprof
Le débit instantané vaut
\[
q(t)=S\dot{\lambda}(t)=eS\dot{\theta}(t)\cos\theta(t).
\]
\fi

\question{En utilisant Python, tracer le débit instantané de la pompe pour un tour de pompe pour $e=\SI{10}{mm}$ et $R=\SI{10}{mm}$ ainsi que pour $e=\SI{20}{mm}$ et $R=\SI{5}{mm}$. La fréquence de rotation est $\dot{\theta}(t)=\SI{100}{rad.s^{-1}}$, la section du piston est donnée par $S=\SI{1}{cm^2}$.}
\ifprof
Pour une vitesse angulaire constante $\omega=\SI{100}{rad.s^{-1}}$, on pose
$\theta(t)=\omega t$. Le débit s'écrit alors
\[
q(t)=eS\omega\cos(\omega t).
\]
Avec $S=\SI{1e-4}{m^2}$, l'amplitude du débit est
\[
q_{\max}=eS\omega.
\]
On obtient donc :
\[
q_{\max}=\SI{1e-4}{m^3.s^{-1}}
\quad\text{pour } e=\SI{10}{mm},
\]
\[
q_{\max}=\SI{2e-4}{m^3.s^{-1}}
\quad\text{pour } e=\SI{20}{mm}.
\]
Le rayon $R$ modifie la position moyenne du piston mais pas l'amplitude du débit.

\begin{marginfigure}
\includegraphics[width=\linewidth]{11_02_c}
\end{marginfigure}
\fi

\ifprof
\else
\footnotesize
\ifcolle
\else
\marginnote{
\begin{solution}
\begin{enumerate}
\item Tracer les solides et les liaisons du mécanisme.
\item $e\sin\theta+R-\lambda(t)=0$.
\item $\dot{\lambda}(t)=e\dot{\theta}(t)\cos\theta(t)$.
\item $q(t)=eS\dot{\theta}(t)\cos\theta(t)$.
\item Pour $\theta(t)=\omega t$, tracer $q(t)=eS\omega\cos(\omega t)$.
\end{enumerate}
\end{solution}}
\fi
\normalsize

\marginnote{Corrigé voir \ref{GEO:03:C2:06:11}.}
\fi

% BEGIN AUTO CORRIGE INCLUDE
\InclureCorrige{corrige.tex}
% END AUTO CORRIGE INCLUDE
